\chapter{Islet-Cell Autoantibody Test}

\section*{What Is This Test?}
Islet-cell autoantibody testing detects antibodies directed against pancreatic islet-cell antigens. These antibodies indicate autoimmune injury to insulin-producing beta cells and can support the diagnosis or classification of type 1 diabetes and slowly progressive autoimmune diabetes in adults.

\section*{Alternative Names}
Islet Cell Antibodies; ICA; Pancreatic Islet-Cell Autoantibodies; Diabetes Autoantibody Testing

\section*{How It Is Performed}
The antibodies are measured in serum. Modern diabetes-autoantibody evaluation may include antibodies to glutamic acid decarboxylase (GAD65), insulinoma-associated antigen-2 (IA-2), zinc transporter 8 (ZnT8), and insulin itself, in addition to classic islet-cell cytoplasmic antibodies. The exact panel varies by laboratory and clinical question.

\section*{Why It Is Ordered}
\begin{itemize}
  \item Distinguishing autoimmune type 1 diabetes from type 2 diabetes or other diabetes types when the clinical presentation is uncertain
  \item Evaluating lean-onset diabetes, rapid loss of insulin production, or diabetes with other autoimmune disease
  \item Assessing selected relatives of people with type 1 diabetes in research or specialist risk-stratification programs
\end{itemize}

\section*{Interpretation and Limitations}
A positive result supports autoimmune diabetes but does not by itself measure current insulin production or predict the exact time of diabetes onset. A negative result does not exclude type 1 diabetes, especially when only one antibody is tested or disease is longstanding. Results should be interpreted with glucose, HbA1c, C-peptide, symptoms, age, and clinical course.
\section*{Specimen and Preparation}
The test generally uses serum from a blood sample. Fasting is usually not required for antibody measurement, although glucose, lipid, or other tests collected at the same visit may have separate instructions. The requested panel matters: GAD65, IA-2, ZnT8, insulin autoantibodies, and classic ICA assays do not have identical performance or clinical uses.

\section*{Risks and Safety}
Physical risk is limited to venipuncture. A positive antibody result does not mean that a person is immediately ill and should not delay treatment when hyperglycemia or diabetic ketoacidosis is suspected. Excessive thirst, frequent urination, vomiting, abdominal pain, deep breathing, confusion, or marked weakness requires urgent assessment.

\section*{Understanding Results and Follow-up}
The report should identify each antibody, its result or titre, the assay cutoff, and whether it is positive, negative, or indeterminate. Multiple positive antibodies generally indicate greater likelihood of autoimmune beta-cell disease than one low-level result, but risk is not the same as diagnosis. Follow-up may include glucose, HbA1c, C-peptide, repeat testing, and diabetes education. Do not change treatment based on antibody results alone.

\section*{Regulatory Status and Authorization Date}
Regulatory status applies to the specific assay, instrument, manufacturer, and intended use rather than to the test name alone. No single approval date can be assigned to this test category. For a commercial assay, verify the current product in the relevant regulator's database: the U.S. Food and Drug Administration (FDA) clearance, approval, or authorization; India's Central Drugs Standard Control Organisation (CDSCO) licence or permission; the European Union In Vitro Diagnostic Regulation (EU IVDR) CE certificate issued through the applicable conformity-assessment route; or the United Kingdom Medicines and Healthcare products Regulatory Agency (MHRA) registration and UKCA/CE status. Laboratory-developed tests may not have an individual FDA, CDSCO, EU IVDR, or MHRA approval date.
\clearpage
